\subsection{Canonical-Orbit Boltzmann--Poisson Solver}
\label{sec:fully-kinetic-bp}

\begin{figure}[htbp]
\centering
\textbf{Physics flow: canonical-orbit Boltzmann--Poisson module}
\par\smallskip
\begin{tikzpicture}[node distance=0.22cm]
\node[physicsflow] (cb1) {Compute the spherical gravitational potential from the current matter distribution};
\node[physicsflow,below=of cb1] (cb2) {Build an orbit atlas that records turning points, travel times, and radial occupancy for each orbit family};
\node[physicsflow,below=of cb2] (cb3) {Store the matter by orbit family and position along the orbit};
\node[physicsflow,below=of cb3] (cb4) {Move that matter forward along the frozen orbits until the next field or collision event};
\node[physicsflow,below=of cb4] (cb5) {At a collision event, sample the local incoming velocities and choose physical collision pairs};
\node[physicsflow,below=of cb5] (cb6) {Scatter accepted pairs and place their outgoing matter onto the corresponding new trajectories};
\node[physicsflow,below=of cb6] (cb7) {Return delayed central trajectories, compact the continuous packet list, and close the mass and energy ledgers};
\node[physicsflow,below=of cb7] (cb8) {Rebuild density, gravity, and the orbit atlas when the field changes; then schedule the next event};
\draw[physicsarrow] (cb1) -- (cb2);
\draw[physicsarrow] (cb2) -- (cb3);
\draw[physicsarrow] (cb3) -- (cb4);
\draw[physicsarrow] (cb4) -- (cb5);
\draw[physicsarrow] (cb5) -- (cb6);
\draw[physicsarrow] (cb6) -- (cb7);
\draw[physicsarrow] (cb7) -- (cb8);
\end{tikzpicture}
\end{figure}

\subsubsection{Purpose and state variables}

This module evolves a spherical SIDM distribution under self-gravity and
finite-angle elastic scattering.  Its state contains two connected
representations of the same physical system:
\begin{enumerate}
\item a finite-volume distribution on an orbit grid, used for long intervals
      of collisionless motion; and
\item continuous phase space packets created by recent collisions, used until
      they can be placed on the orbit grid with a conservative map.
\end{enumerate}

The implementation is organized in
\path{cpu/include/ev/canonical_orbit_phase_transport.hpp},
\path{cpu/include/ev/canonical_boltzmann_poisson_event.hpp}, and
\path{cpu/include/ev/weighted_vlasov_pic_collisions.hpp}.  Restarts are handled by
\path{cpu/include/ev/canonical_boltzmann_poisson_restart.hpp}.

In spherical symmetry we use
\[
  \eta=\frac{L^2}{2},\qquad
  H=\Phi(r)+\frac{w^2}{2}+\frac{\eta}{r^2}.             \tag{CBP1}
\]
Here \(w=v_r\), \(L=r v_t\), and \(H\) is the specific Hamiltonian in the
current gravitational potential.  The phase space measure is
\[
  \dd\Gamma=8\pi^2\,\dd r\,\dd w\,\dd\eta.             \tag{CBP2}
\]
The factor \(8\pi^2\) is included in every physical mass budget.

For a frozen spherical potential, \(\eta\) and \(H\) remain constant along an
orbit.  A third coordinate, the dimensionless radial phase
\(\zeta_{\rm orb}\in[0,1)\), tells us where the particle is between
pericentre and apocentre.  One radial period advances \(\zeta_{\rm orb}\) by
one.  The subscript distinguishes this orbit phase from the Schwarzschild time
\(\tau\) used in the relativistic chapters.  The orbit-grid coordinates are
therefore
\[
  (\eta,H,\zeta_{\rm orb}).                             \tag{CBP3}
\]

\noindent\textbf{Module interface.}
A dedicated caller initializes the canonical-orbit library.  The generic
\path{params.ini} interface selects the Eulerian methods described elsewhere.
The canonical initializer takes a cell-integrated distribution, an equilibrium
table \(f(H)\), the resolution of the orbit atlas, collision-shell controls,
field and collision clocks, and an output directory.  A continuation additionally reads the canonical restart,
which stores the atlas, continuous packets, random counters, and both event
clocks.  The task-specific Eulerian finite-angle frontends described in
Section~\ref{sec:eulerian-grid-collisions} are separate algorithms and must
not be used as substitutes for this orbit-atlas representation.

\subsubsection{Building the canonical orbit atlas}

The orbit atlas tabulates the bound trajectories represented by the
finite-volume cells in equation~(CBP3).  Its construction first fixes angular
momentum and energy, then resolves orbital phase and the overlap with the
collision grid.  These four steps are described below.

\paragraph{Step 1: angular momentum planes.}

Each Eulerian \(\eta\) cell is integrated with a small Gaussian rule.  The
first cell uses a left Gauss--Radau rule, so one plane lies exactly at
\(\eta=0\).  The last cell uses a right Gauss--Radau rule.  Interior cells use
Gauss--Legendre points.  These endpoint planes give direct support to radial
orbits and to the largest represented angular momentum.

The rule selection appears in the atlas builder:
\codesource{Canonical orbit-atlas angular momentum quadrature}
{\path{cpu/include/ev/canonical_orbit_phase_transport.hpp};
 \path{cpu/include/ev/phase_space_cells.hpp}}
{Chooses endpoint-inclusive Gauss--Radau planes at the two angular momentum
 boundaries and Gauss--Legendre planes in interior cells}
\begin{lstlisting}[language=C++]
if (eta_cell == 0)
  return phase_space_gauss_radau_left_rule(points);
if (eta_cell == p.NL - 1)
  return phase_space_gauss_radau_right_rule(points);
return phase_space_gauss_legendre_rule(points);
\end{lstlisting}

\paragraph{Step 2: energy rows.}

At fixed \(\eta\), define the effective potential
\[
  \Phi_{\rm eff}(r;\eta)=\Phi(r)+\frac{\eta}{r^2}.       \tag{CBP4}
\]
For a bound orbit with energy \(H\), the turning points solve
\(\Phi_{\rm eff}=H\).  The inner turning radius
\(r_p\) is the pericentre, and the outer turning radius \(r_a\) is the
apocentre; radial motion is allowed for \(r_p\le r\le r_a\).  The radial
action is
\[
  J_r(H,\eta)=\frac{1}{\pi}
  \int_{r_p}^{r_a}\sqrt{2[H-\Phi_{\rm eff}(r;\eta)]}\,\dd r.  \tag{CBP5}
\]
Energy-row faces are placed through monotone edges in \(J_r\).  An optional
positive weight based on the initial \(f(H,\eta)\) concentrates rows where
the phase space mass is large.  The row stores both its action interval and
its energy interval.  During a later Poisson rebuild, the action edges define
the corresponding energy edges in the new potential.

\paragraph{Step 3: phase cells.}

For one representative orbit, let \(T_r\) be the radial period.  The orbit
phase is elapsed orbital time normalized by this period,
\[
  \zeta_{\rm orb}=\frac{t_{\rm orbit}}{T_r}\pmod 1,
  \qquad \dot\zeta_{\rm orb}=\frac{1}{T_r}.             \tag{CBP6}
\]
Each energy row is divided into equal intervals in \(\zeta_{\rm orb}\).
Equal phase width means equal residence-time width.  Dense sampling near a
turning point then arises from the orbit map itself, where a fixed time
interval covers a small radial interval.

\paragraph{Step 4: overlap with the collision grid.}

Density, pressure, collisions, and the Poisson solve use the Eulerian cells
\((r,w,\eta)\).  The atlas builder traces the boundary of each canonical cell
into the \((r,w)\) plane and intersects it with the Eulerian rectangles.  For
each intersection it stores
\[
  \int\dd\Gamma,\qquad
  \int H\,\dd\Gamma,\qquad
  \int H^2\,\dd\Gamma,                                  \tag{CBP7}
\]
together with the geometric moments needed by the quadratic reconstruction.
These overlaps are computed once for one frozen potential and reused by all
projections in that interval.

\begin{figure}[ht]
\centering
\begin{tikzpicture}[x=1.0cm,y=1.0cm,>=Latex]
  \draw[step=1cm,gray!55,thin] (0,0) grid (6,4);
  \node[below] at (3,0) {$r$};
  \node[rotate=90] at (-0.45,2) {$w$};
  \fill[blue!18]
    (0.65,0.65) .. controls (1.45,1.10) and (2.25,1.55) ..
    (3.20,2.05) .. controls (4.10,2.55) and (5.05,3.05) ..
    (5.55,3.45) --
    (5.15,3.85) .. controls (4.35,3.28) and (3.55,2.90) ..
    (2.75,2.50) .. controls (1.80,2.02) and (1.00,1.52) ..
    (0.35,1.05) -- cycle;
  \draw[blue!70!black,very thick]
    (0.65,0.65) .. controls (1.45,1.10) and (2.25,1.55) ..
    (3.20,2.05) .. controls (4.10,2.55) and (5.05,3.05) ..
    (5.55,3.45);
  \draw[blue!70!black,very thick]
    (0.35,1.05) .. controls (1.00,1.52) and (1.80,2.02) ..
    (2.75,2.50) .. controls (3.55,2.90) and (4.35,3.28) ..
    (5.15,3.85);
  \node[fill=white,draw,rounded corners=2pt,align=center,inner sep=2pt]
    at (1.45,3.45)
    {one finite cell in\\$ (H,\zeta_{\rm orb}) $};
  \draw[->,thick] (1.90,3.10) -- (2.85,2.20);
  \node[align=center,fill=white,inner sep=2pt] at (4.55,0.38)
    {stored overlap moments};
  \draw[->,thick] (4.55,0.68) -- (4.25,2.67);
\end{tikzpicture}
\caption{One canonical cell mapped into the Eulerian \((r,w)\) plane at fixed
\(\eta\).  Its curved image crosses several rectangles.  VENUS stores the
integral over each shaded intersection, so later projections use integrated
mass instead of a single sample point.}
\label{fig:canonical-atlas-overlap}
\end{figure}

Each canonical phase cell stores three integrated quantities,
\[
  M_c=\int_c f\,\dd\Gamma,\qquad
  E_c=\int_c Hf\,\dd\Gamma,\qquad
  E_{2,c}=\int_c H^2f\,\dd\Gamma.                        \tag{CBP8}
\]
They retain the mean and variance of \(H\) inside an energy row.  A quadratic
reconstruction in \(H\) can therefore be projected without replacing the row
by its center energy.  Positivity is checked from the reconstructed value;
the stored cell moments are also checked against the exact moment bounds for
the interval \([H_L,H_R]\):
\[
 H_L\le \langle H\rangle\le H_R,\qquad
 \langle H\rangle^2\le\langle H^2\rangle
 \le(H_L+H_R)\langle H\rangle-H_LH_R.                  \tag{CBP9}
\]

\subsubsection{Frozen-potential Vlasov transport}

During a frozen-field interval, each row advances by a phase shift,
\[
  \zeta_{\rm orb}\longmapsto
  \zeta_{\rm orb}+\frac{\Delta t}{T_r}\pmod 1.         \tag{CBP10}
\]
The finite-volume masses and Hamiltonian moments stay in their canonical
cells.  The row stores the accumulated shift in \code{phase_offset}:
\codesource{Canonical frozen-orbit phase transport}
{\path{cpu/include/ev/canonical_orbit_phase_transport.hpp}}
{Advances each orbit row by its period-normalized phase translation and wraps
 the result onto the unit radial-phase interval}
\begin{lstlisting}[language=C++]
row.phase_offset = hamiltonian_orbit_wrap_phase(
    row.phase_offset + dt / row.period);
\end{lstlisting}

This update is an exact translation for the representative orbit in the
frozen potential.  Projection to \((r,w,\eta)\) is performed when moments or
collisions are required.  Consequently many orbital intervals can be
accumulated before a new reconstruction is needed.

Nearly radial trajectories use the regularized orbit geometry in
\path{cpu/include/ev/hamiltonian_orbit_phase.hpp}.  The ordinary integrator is
used when its energy check passes.  A failed check triggers the
Levi--Civita-style regularized evaluation near pericentre.  The orbit table
records retry and unresolved-orbit diagnostics.

\subsubsection{The finite inner radius}
\label{sec:delayed-inner-carrier}

The logarithmic radial mesh begins at \(R=r_{\min}>0\).  An orbit that enters
\(0<r<R\) keeps evolving in the central potential.  Its canonical phase map
contains a hidden phase interval whose duration is the physical inward and
outward flight time.  Mass in this phase interval is included in the enclosed
mass used by Poisson.  When the phase reaches the outgoing crossing, the same
orbit returns through \(R\) with outward radial velocity.

The resolved and hidden pieces obey
\[
 M_{\rm atlas}=M_{\rm visible}+M_{\rm hidden}+M_{\rm escaped}. \tag{CBP11}
\]
The code checks this partition for every canonical column.  SIDM calculations
also report the optical depth across \(0<r<R\) and are repeated with smaller
\(R\).

\subsubsection{Choosing the next event time}

The potential and collision clocks are independent.  From consecutive
Poisson fields, VENUS estimates a mass-weighted field-change time using both
\(\Phi\) and \(g=\partial_r\Phi\).  The requested field interval is a fraction
\(\epsilon_\Phi\) of that time.  The collision module supplies a physical
collision time and an NTC-majorant time.  The scheduler uses
\[
 \Delta t=\min\left(
  \Delta t_{\max},
  \epsilon_\Phi t_\Phi-t_{\Phi,{\rm age}},
  \epsilon_C t_C-t_{C,{\rm age}},
  \epsilon_{\rm maj}t_{\rm maj}
 \right).                                               \tag{CBP12}
\]
Here \(\Delta t_{\max}\) is the configured hard cap;
\(t_\Phi\) is the mass-weighted field-change timescale, and
\(t_{\Phi,{\rm age}}\) is the elapsed time since the last potential rebuild.
The derivative estimate is the smaller of
\(\sqrt{\sum_i\omega_i\Phi_i^2/\sum_i\omega_i\dot\Phi_i^2}\) and
\(\sqrt{\sum_i\omega_i g_i^2/\sum_i\omega_i\dot g_i^2}\), where
\(g=\Phi_r\), \(M_i=4\pi\max(\rho_i,0)
(r_{i+1/2}^3-r_{i-1/2}^3)/3\) is nonnegative resolved shell mass, and
\(\omega_i=M_i/(M_{\rm hidden}+\sum_jM_j)\).
Here \(M_{\rm hidden}\) is the unresolved inner mass in (CBP11), and the
rates follow from \code{vp_potential_time_derivative}.  Each channel requires
both mean squares to exceed \((4096\epsilon_{\rm mach})^2\), with
\(\epsilon_{\rm mach}\) double-precision machine roundoff. A channel at or
below this threshold gives infinite time, as do zero total mass and the
helper's expanding-background branch.
After an event the solver also restricts \(t_\Phi\) by the elapsed field time
divided by the larger mass-weighted relative potential or force change
returned by \code{canonical_field_mismatch}.  These are the distinct
canonical helpers in \path{cpu/include/ev/canonical_boltzmann_poisson_event.hpp},
lines 242--418 and 1801--1835.
Likewise \(t_C\) is the minimum local physical
collision timescale and \(t_{C,{\rm age}}\) is the elapsed time since the
last collision event.  The dimensionless fractions \(\epsilon_\Phi\) and
\(\epsilon_C\) set the corresponding deadlines.  Finally,
\(t_{\rm maj}\) is the no-time-counter (NTC) proposal-majorant timescale and
\(\epsilon_{\rm maj}\) is its safety fraction.  Subtracting each age prevents
one kind of event from indefinitely postponing the other.

The source call is:
\codesource{Canonical Boltzmann--Poisson event scheduler}
{\path{cpu/include/ev/canonical_boltzmann_poisson_event.hpp};
 \path{cpu/include/ev/weighted_vlasov_pic_event_scheduler.hpp}}
{Combines the independent field and collision clock ages with their physical
 deadlines to choose the next event interval and ordering}
\begin{lstlisting}[language=C++]
diagnostics.plan = plan_weighted_pic_event(
    {state.potential_age, state.collision_age},
    state.potential_timescale,
    scheduled_collision_timescale,
    options.potential_fraction,
    options.collision_fraction,
    maximum_interval);
\end{lstlisting}

If the collision clock expires first, the sequence is orbit advance,
collision, and later field rebuild.  If the field clock expires first, the
sequence is orbit advance, field rebuild, and later collision.  When both
expire together, the scheduler records their chronological order and executes
both operations at the same event time.

\subsubsection{Sampling the instantaneous distribution}

At a collision event, the atlas and all active continuous packets are sampled
at their current phase.  A sampled marker contains
\[
  (r,w,v_{t,1},v_{t,2},m_p).                             \tag{CBP13}
\]
The tangential direction is uniform in its local plane.  Marker weights are
fractional, so a canonical cell can contribute to several independent virtual
batches.  Increasing the number of markers improves the estimate of the local
collision integral without changing the orbit-grid resolution.

Markers are sorted by radius and divided into disjoint shells.  A shell is
accepted for local scattering when
\[
  \Delta r_s\le\lambda_{\rm rms,s}.                     \tag{CBP14}
\]
The shell builder can increase its population to reduce statistical noise.
The ledger records the shell width, mean free path, marker count, and the
largest ratio \(\Delta r_s/\lambda_{\rm rms,s}\).

\subsubsection{No-time-counter pair selection}

Consider one shell of volume \(V_s\), density \(\rho_s\), and \(N_s\) sampled
markers.  Define
\[
  A_s\ge \max_{i,j}\left[u_{ij}\sigma_{\rm tot}(u_{ij})\right],
  \qquad u_{ij}=|\boldsymbol v_i-\boldsymbol v_j|.       \tag{CBP15}
\]
The no-time-counter clock proposes unordered pairs at rate
\[
  \nu_{{\rm trial},s}=\frac{N_s-1}{2}\rho_s A_s.        \tag{CBP16}
\]
For every proposal, two distinct marker indices are drawn uniformly.  With a
fresh uniform random number \(U_a\in(0,1)\), the pair is accepted when
\[
  U_a A_s<u_{ij}\sigma_{\rm tot}(u_{ij}).               \tag{CBP17}
\]
Equations~(CBP16)--(CBP17) give the required rate
\(f_i f_j u_{ij}\sigma_{\rm tot}\) in expectation.  The algorithm visits the
realized proposals and accepted events.  Its cost scales with the physical
trial count.

The implementation follows these equations directly:
\codesource{Canonical-orbit no-time-counter sampler}
{\path{cpu/include/ev/canonical_orbit_phase_transport.hpp}}
{Draws distinct local phase space samples and accepts each trial against the
 shell's verified relative-speed--cross-section majorant}
\begin{lstlisting}[language=C++]
clock += -std::log(rng.uniform_open()) / trial_rate;
ia = rng.uniform_index(count);
ib = draw_distinct_index(count, ia);
u_sigma = relative_speed * total_cross_section_per_mass(...);
if (rng.uniform_open() * shell.majorant_u_sigma >= u_sigma)
  continue;
\end{lstlisting}

\subsubsection{Drawing the outgoing direction}

Let
\[
 \boldsymbol V=\frac{\boldsymbol v_1+\boldsymbol v_2}{2},
 \qquad \boldsymbol u=\boldsymbol v_1-\boldsymbol v_2. \tag{CBP18}
\]
Elastic scattering changes the direction of \(\boldsymbol u\) and keeps its
magnitude.  Write \(\mu=\cos\theta\).  The normalized cumulative distribution
for \(\mu\) is
\[
 P(\mu\le x)=
 \frac{2\pi}{\sigma_{\rm tot}(u)}
 \int_{-1}^{x}\frac{\dd\sigma}{\dd\Omega}(u,\mu')\,\dd\mu'. \tag{CBP19}
\]
VENUS draws \(U_\mu\in(0,1)\) and inverts equation~(CBP19).  Isotropic
scattering gives \(\mu=2U_\mu-1\).  For the Born Yukawa kernel, the analytic
inverse used by the code is
\[
 a=m_\phi^2,\qquad b=\frac{m_\chi^2u^2}{2},
\]
\[
 y=\frac{2b}{a},\qquad
 \mu=1-\frac{2U_\mu}{1+y(1-U_\mu)}.                     \tag{CBP20}
\]
Here \(U_\mu\) samples the probability of a cosine at or above \(\mu\),
complementary to the increasing cumulative probability in (CBP19).
This factored expression has the continuous uniform-cosine limit
\(\mu=1-2U_\mu\) at \(y=0\); the numerical inverse uses \(1/y\) when
\(y>1\) to keep the arithmetic bounded.
A second independent random number gives the azimuth
\(\psi=2\pi U_\psi\).  The code constructs the unit vector
\(\widehat{\boldsymbol u}'(\mu,\psi)\) around the incoming relative direction
and sets
\[
 \boldsymbol v'_1=\boldsymbol V+\frac{u}{2}
       \widehat{\boldsymbol u}',\qquad
 \boldsymbol v'_2=\boldsymbol V-\frac{u}{2}
       \widehat{\boldsymbol u}'.                        \tag{CBP21}
\]
Thus each accepted pair satisfies
\[
 \boldsymbol v'_1+\boldsymbol v'_2=\boldsymbol v_1+\boldsymbol v_2,
 \qquad
 |\boldsymbol v'_1|^2+|\boldsymbol v'_2|^2
 =|\boldsymbol v_1|^2+|\boldsymbol v_2|^2.              \tag{CBP22}
\]

The corresponding source excerpt is:
\codesource{Finite-angle elastic-scattering map}
{\path{cpu/include/ev/weighted_vlasov_pic_collisions.hpp};
 \path{cpu/include/ev/finite_angle_boltzmann.hpp}}
{Samples the scattering cosine and azimuth and constructs the two outgoing
 centre-of-mass-frame velocities while preserving pair invariants}
\begin{lstlisting}[language=C++]
const double mu = sample_scatter_mu(u, p, options, rng);
const double psi = 2.0 * PI * rng.uniform_open();
const BPFAVec3 direction =
    bp_fa_scattered_relative_direction(relative, mu, psi);
const BPFAVec3 relative_after = bp_fa_scale(direction, u);
const BPFAVec3 out1 = bp_fa_add(centre,
    bp_fa_scale(relative_after, 0.5));
const BPFAVec3 out2 = bp_fa_sub(centre,
    bp_fa_scale(relative_after, 0.5));
\end{lstlisting}

Both tangential components enter equation~(CBP21).  After scattering, the
new angular momentum of each packet is
\[
 L'=r\sqrt{(v'_{t,1})^2+(v'_{t,2})^2},\qquad
 \eta'=\frac{(L')^2}{2}.                                \tag{CBP23}
\]
Collisions therefore redistribute angular momentum while preserving the full
three-dimensional pair momentum and kinetic energy.

\subsubsection{Gain, loss, and reverse events}

An accepted event contains its incoming and outgoing velocities in one
record.  The incoming mass is withdrawn from the canonical cells that owned
the sampled marker.  The two outgoing packets are then created with the same
event mass.  This is the sampled gain-minus-loss update of the Boltzmann
operator.

The collision map in equation~(CBP21) is one-to-one and preserves velocity
volume.  Drawing \(\psi\) over the full interval and drawing \(\mu\) from the
complete differential cross section includes the reverse map with its
physical probability.  For a Maxwellian \(f_M\propto\exp(-v^2/2T)\), equation
(CBP22) gives
\[
 f_M(\boldsymbol v_1)f_M(\boldsymbol v_2)
 =f_M(\boldsymbol v'_1)f_M(\boldsymbol v'_2).            \tag{CBP24}
\]
The expected collision update therefore vanishes.  A numerical realization
should approach this null state as the collision-marker population and orbit
atlas are refined.

\subsubsection{Source-owned withdrawal and outgoing carriers}

Each sampled canonical marker keeps the finite-volume cells and weights that
created it.  Once a pair is accepted, the incoming event mass is removed from
those owned cells immediately.  The source cells move with the frozen-field
phase flow, so this withdrawal includes their resolved and hidden orbital
phases consistently.

The outgoing velocities define continuous packets at the collision radius.
They are advanced along their new orbits until the next collision event.  The
main source sequence is:
\codesource{Canonical source withdrawal and outgoing carriers}
{\path{cpu/include/ev/canonical_orbit_phase_transport.hpp}}
{Removes accepted incoming mass from its owned atlas cells and retains the
 two outgoing continuous packets for conservative remapping onto new orbits}
\begin{lstlisting}[language=C++]
represented_first_before = withdraw_source(
    pair.first_ref, pair.first_source,
    event.first_before, packet_mass);
represented_second_before = withdraw_source(
    pair.second_ref, pair.second_source,
    event.second_before, packet_mass);

retain_signed_marker(new_signed_positive,
    collision_source[pair.first_source],
    event.first_after, packet_mass);
retain_signed_marker(new_signed_positive,
    collision_source[pair.second_source],
    event.second_after, packet_mass);
\end{lstlisting}

At creation, the atlas loss and outgoing packet gain satisfy
\[
 \Delta M_{\rm atlas}+M_{\rm outgoing}=0.              \tag{CBP25}
\]
The pair ledger also checks equation~(CBP22) for every accepted event.

\subsubsection{Mapping outgoing packets onto the orbit grid}

Before the next collision distribution is sampled, active outgoing packets
are placed on the current atlas:
\begin{enumerate}
\item Compute \(\eta'\) and \(H'\) from equations~(CBP1) and (CBP23).
\item Find the two neighboring \(\eta\) quadrature planes with positive linear
      weights.
\item Find the energy row containing \(H'\) on each plane.
\item Integrate the orbit to obtain the physical phase
      \(\zeta_{\rm orb}'\).
\item Convert \(\zeta_{\rm orb}'\) to the co-moving phase coordinate using the row's
      current \code{phase_offset}.
\item Deposit mass, \(H\)-mass, and \(H^2\)-mass with positive weights into at
      most four neighboring canonical cells.
\end{enumerate}

The deposition weights sum to one.  The materialization transaction checks
mass before committing it.  It also checks the moment bounds in
equation~(CBP9).

An outgoing packet with \(H'>0\) lies above the bound-orbit atlas.  It remains
an explicit open-orbit packet.  The characteristic integrator advances it,
the collision sampler includes it while it is inside the radial domain, and
the escape ledger records it after it crosses the outer boundary.  Separate
counters classify open packets caused by high energy, angular momentum
support, a missing row, an energy gap, or an orbit-geometry failure.

\subsubsection{Keeping the continuous packet list compact}

Many small collision packets may occupy nearly the same local phase space
region.  VENUS replaces each such group by a positive quadrature supported on
a subset of the existing packet locations.  Positions and velocities remain
physical packet states.  The new nonnegative weights preserve fourteen group
moments:
\begin{align*}
&1, q, w, qw, w^2, qw^2, v_t^2, qv_t^2,\\
&\frac{v^2}{2}w, q\frac{v^2}{2}w,
 \eta, q\eta, H, H^2,\qquad q=\frac{r^3}{3}.          \tag{CBP26}
\end{align*}
These constraints preserve mass, radial momentum, radial and tangential
kinetic moments, radial kinetic-energy flux, the squared-angular momentum
moment \(\eta=L^2/2\), and the first
two Hamiltonian moments.  The factors with \(q\) preserve the same quantities
under radial cloud-in-cell deposition.

The reduction is local in \((r,w,\eta)\).  A group contains packets with the
same deposition stencil.  A Carath\'eodory reduction is a positive
convex-geometry compression: it replaces a nonnegative weighted cloud by a
subset of its existing packet locations while preserving the requested
moment vector.  It leaves at most fourteen packets in a well-conditioned
group.  Here ``well-conditioned'' means that the retained moment columns are
sufficiently independent for their nonnegative weights to be solved without
amplifying roundoff.  Every group recomputes equation~(CBP26) after reduction.
A group whose moment matrix is ill-conditioned or whose residual exceeds the
requested tolerance keeps its original packets.

Packet ordering is made deterministic from marker identity and phase space
coordinates before the reduction, so a fixed platform and random seed retain
the same packet order across supported OpenMP thread counts.

\subsubsection{Refreshing the gravitational field}

When the field clock expires, VENUS projects the instantaneous atlas and
continuous packets to radial moments and solves spherical Poisson.  A new
orbit atlas is built in the new potential with the previous radial-action
edges.  Mass and its \((H,H^2)\) moments are transferred between matching
topological rows.  The field iteration stops when the mass-weighted mismatch
between the represented field and a fresh Poisson solve is below tolerance.

This rebuild lets the slowly changing stationary part of \(f(H,\eta)\) evolve
with the potential.  The current distribution supplies the transferred row
moments; the initial equilibrium table is used only for initialization.

\subsubsection{Conservation diagnostics and output}

Every event writes the following independent checks:
\begin{itemize}
\item atlas mass before and after frozen transport;
\item represented mass before and after collisions;
\item pair momentum and kinetic-energy errors;
\item shell locality \(\Delta r_s/\lambda_{\rm rms,s}\);
\item materialization residuals in \(M\), \(H\), and \(H^2\);
\item packet-compression moment residuals;
\item active, hidden, and escaped open-orbit mass;
\item field mismatch and the measured potential-change time;
\item marker-batch estimates of density and radial heat-flux noise.
\end{itemize}

The main files are \path{history.csv}, \path{profiles.csv}, and
\path{operator_mass_budget.csv}.  A restart stores the atlas, phase offsets,
all three cell moments, open packets, deferred outgoing packets, random-stream
counters, and both event clocks.

Convergence studies should vary the orbit-grid resolution, marker count,
virtual collision batches, local shell population, compression subcells,
\(r_{\min}\), and the allowed active open-tail fraction.  The comparison uses
density, radial velocity dispersion, collision energy transfer, and radial
heat flux.  A gravothermal core-collapse comparison should additionally
resolve the early dilution, subsequent reversal, and central-density growth
through the chosen physical interval.
